% RESEARCH_DRAFT_SHA256: c58e4e7be7c530cf7d9167979d3c66bc98b7d75d27c2613263b686b836ea8e7f
% Layout and source-reference wrappers expanded for the web renderer.
\documentclass[11pt]{article}
\usepackage[margin=25mm]{geometry}
\usepackage{amsthm}
\usepackage{fontspec,unicode-math}
\setmainfont{Latin Modern Roman}
\setmathfont{Latin Modern Math}
\usepackage{amsmath,amssymb,mathtools}
\usepackage{xurl,hyperref}
\hypersetup{colorlinks=true,urlcolor=blue,linkcolor=blue}
\setcounter{secnumdepth}{0}
\setlength{\parindent}{0pt}
\setlength{\parskip}{5pt}
\setlength{\emergencystretch}{3em}
\newcommand{\sref}[2]{\hyperlink{source:#1:#2}{[S#2]}}
\newcommand{\cataloguescope}[1]{\paragraph{Source question.}#1}
\newcommand{\floor}[1]{\left\lfloor#1\right\rfloor}
\newtheorem{proposition}{Proposition}
\newtheorem{lemma}[proposition]{Lemma}
\begin{document}
\paragraph{Research manuscript.} OpenAI \texttt{gpt-6-astra}, reasoning \texttt{ultra}; 11 October 2026.
\paragraph{Preserved source body.} The complete catalogue mathematical text and sources follow with layout and source links expanded, including source qualifications and its historical preparation notes. New research begins at the section headed Research attempt.

% UnsolvedMath ID 30004009; source catalogue code OWR-16634-002
\section{Real Tau Conjecture for Depth-Four Circuits}\label{30004009}
{\small UnsolvedMath ID 30004009 · Computer Science}
\subsection{Definitions and mathematical statement}

\paragraph{Shared notation.}Standard number systems are denoted by $\mathbb N,\mathbb Z,\mathbb Q,\mathbb R,\mathbb C$; index ranges, computational models and specialized notation follow the cited problem source and the definitions below.


A sparse univariate polynomial has few nonzero monomials, with exponents encoded as prescribed. The stated depth-four form sums products of such sparse polynomials; real roots are counted without multiplicity.

An arithmetic circuit is a directed acyclic graph whose input nodes are variables or field constants and whose internal nodes perform additions and multiplications (with subtraction allowed when stated). Size counts gates; depth is the maximum number of gates on an input-output path.

A Boolean function maps $\{0,1\}^n$ to $\{0,1\}$. A Boolean circuit is an acyclic network of logical gates; a formula is a circuit with no sharing of intermediate gate outputs. Size and depth measure its number of gates and longest gate path.
\paragraph{Statement recorded in the supplied source.}
Is the number of distinct real zeros of every univariate real polynomial represented by a depth-four arithmetic circuit, consisting of a sum of products of sparse polynomials, bounded by a polynomial in the circuit size? \href{https://doi.org/10.4171/owr/2018/51}{[S2]}
\paragraph{Scope qualification.}
The size convention must account for the source\textquotesingle s sparsity, products, and exponent encoding; no degree bound is inserted.
\paragraph{Source status.}
The supplied catalogue labels this record partially solved. Its source-status metadata was checked against the supplied snapshot on 2026-10-11. This import does not independently establish current open status or mathematical correctness. The catalogue literature assessment is dated 2026-08-21.
\subsection{Short English statement}
Bound the number of distinct real roots by a polynomial in the displayed arithmetic representation size.
\subsection{Sources}
\begin{itemize}
\item[S1] Ulam AI and the UnsolvedMath Contributors, UnsolvedMath: A Curated Collection of Open Mathematics Problems, record 30004009 (OWR-16634-002). Catalogue curation and metadata: CC BY 4.0. The supplied source question is retained; definitions and the English summary are editorial additions.
\url{https://huggingface.co/datasets/ulamai/UnsolvedMath}.
\item[S2] Complexity Theory (2019), Oberwolfach Reports, DOI 10.4171/owr/2018/51, pp. 3025--3080
\url{https://doi.org/10.4171/owr/2018/51}.
\item[S3] The Real Tau-Conjecture Is True on Average
\url{https://arxiv.org/abs/1806.00417}.
\item[S4] Oberwolfach Report 51/2018
\url{https://ems.press/content/serial-article-files/46773}.
\textit{Underlying Oberwolfach Reports material retains its source-specific licence, including CC BY-SA 4.0 where applicable; see THIRD\_PARTY\_NOTICES.md.}
\end{itemize}
\paragraph{Preparation and review.}
Original question attributed to the cited source and UnsolvedMath contributors. Codex (GPT-6; exact serving revision not exposed) assisted with definitions, English summaries, source inspection and formatting. No new human mathematical authorship or independent proof verification is asserted. Canonical ID 30004009 and source code OWR-16634-002 checked against the complete supplied registry and existing repository statements; known flagged repetitions and obvious identical conjectures excluded. Full mathematical-equivalence/literature review has not been performed. Statement preparation and final review followed contribution prompts C and D; automated checks verify structure and identity, not open status.
\label{end:30004009}

\section{Research attempt}
\subsection{Literature, source context and scope}
The full original catalogue body and source qualifications are preserved above. The relevant primary workshop passages were read in the cached publisher report. Historical claims and subsequent attributed results are distinguished from the mathematical work below. A restricted example, standard lemma or known theorem is not counted as new progress on the entire original target. Missing assumptions and repeated or rhetorical extracts are not silently turned into different solved problems.

\subsection{The product case works; summation is the missing step}
The primary report and the corrected 2019 version of Briquel--Bürgisser distinguish an average-case theorem from the unrestricted real tau question. We test a direct root-counting route. We explicitly exclude the zero polynomial: it has infinitely many zeros and the usual finite-root conjecture concerns nonzero polynomials, although the short preserved extract does not spell out this convention. All exponents below are nonnegative integers; no bound on their magnitude is assumed.
\begin{proposition}
A nonzero real univariate polynomial with $t$ nonzero monomials has at most $t-1$ distinct positive roots and at most $t-1$ distinct negative roots. A nonzero product $F=\prod_{j=1}^m f_j$ with $t_j$ nonzero monomials in $f_j$ has at most $2\sum_j(t_j-1)+1$ distinct real roots.
\end{proposition}
\begin{proof}
For positive roots, induct on $t$. A single monomial has none. For $t>1$, divide by the smallest power of $x$ appearing; this changes no positive root and leaves a nonzero constant term. The derivative has exactly $t-1$ nonzero monomials, so by induction has at most $t-2$ positive roots. If the original quotient had $r$ distinct positive roots, Rolle's theorem would give at least $r-1$ distinct positive derivative roots, hence $r\le t-1$. Apply the same argument to $f(-x)$ for negative roots. A product vanishes precisely where at least one factor vanishes. Summing the positive and negative bounds over factors and allowing the single possible root zero proves the product bound.
\end{proof}
The same union-of-roots argument is invalid for a sum. Neither $x^2+1$ nor $-2x$ has a positive root, yet their sum $(x-1)^2$ vanishes at 1. More generally, expanding a sum of $k$ products of $m$ factors with at most $t$ monomials gives at most $kt^m$ monomials, so the elementary bound obtained after expansion is at most $2kt^m-1$. This is exponential in the product length $m$, not the requested polynomial dependence on the displayed representation size. Cancellations may reduce it, but this route supplies no suitable bound on them.
\subsection{Verification and remaining gaps}
The cancellation example was evaluated exactly, and sparse products with prescribed integer roots were expanded to compare their distinct roots with the product bound. The zero-polynomial convention, distinct-root count and lack of a degree restriction were audited. No deterministic bound for arbitrary sums of sparse products is established. The known average-case result is attributed and does not justify a worst-case conclusion; no original-target completion is claimed.

\section{References and source attribution}
Primary workshop report: \url{https://doi.org/10.4171/owr/2018/51}. Relevant passages read: 2018/51 p. 3032, real tau conjecture and the separate Gaussian average-case result.

\url{https://arxiv.org/abs/1806.00417}. Read primary corrected-v2 abstract: average-case expected root bound with independent Gaussian coefficients; the page notes correction of an error in v1. This does not settle the unrestricted conjecture.

All retrievals were checked on 11 October 2026. Original question attribution remains in the preserved Sources section. External theorems are attributed and not claimed as new discoveries or fully verified merely by citation.


\paragraph{Integration review.} Codex (GPT-6; exact serving revision not exposed) reviewed the written arguments and expanded layout and source-reference wrappers. Repository metadata and the production web renderer were checked using Python 3.13 and JavaScript with MathJax 3.2.2. This model-assisted review is not independent human proof verification; no PDF compilation or formal proof checking is asserted.

\section{Completion Estimate}
\noindent\textbf{COMPLETION: 0\%}
No substantiated new progress toward the entire original target is claimed. The worked results and failed route are delimited above; a known result, illustrative model, or auxiliary obstruction is not counted as original-target completion.

\section{Final status and attribution}
\textbf{Attempt status: unresolved.} This describes this research attempt, not a new ruling on the source catalogue's status or on external literature claims.
The mathematics and this self-audit were generated by OpenAI \texttt{gpt-6-astra}, reasoning \texttt{ultra}, on 11 October 2026. Codex (GPT-6; exact serving revision not exposed) orchestrated the batch, reviewed the written mathematical arguments, and contributed formatting/repository review. This additional model review is not independent human verification. No human mathematical authorship, independent human proof checking, or formal verification is claimed. The historical preparation paragraphs in the reproduced catalogue body concern that body only.

Contribution prompt A was applied to the complete definition and supplied source record, followed by the adversarial checks of prompt B. Local reproduction files are \texttt{batch022/checks.py}, \texttt{computations.json}, and the per-ID audit JSON. They were not public links at preparation time. Numerical checks are finite and do not prove unrestricted results. No LaTeX engine was found locally; compilation was not performed by this research agent.

\end{document}
